\section*{Capítulo \ref{ch:g10:concentration-and-dilution} --- Concentração e diluição}

\begin{solution}{exo:g10:concentration-and-dilution:1}
$C_m = \qty{5.0}{g} / \qty{0.250}{L} = \qty{20}{g/L}$.
\end{solution}

\begin{solution}{exo:g10:concentration-and-dilution:2}
$c = \qty{0.050}{mol} / \qty{0.200}{L} = \qty{0.25}{mol/L}$.
\end{solution}

\begin{solution}{exo:g10:concentration-and-dilution:3}
$M(\ce{CuSO4}) = 63.5 + 32.1 + 4 \times 16.0 = \qty{159.6}{g/mol}$;
$n = 0.10 \times 0.1000 = \qty{0.0100}{mol}$; $m = 0.0100 \times 159.6 =
\qty{1.60}{g}$.
\end{solution}

\begin{solution}{exo:g10:concentration-and-dilution:4}
$F = 1.0 / 0.050 = 20$; $V_0 = \qty{100.0}{mL} / 20 = \qty{5.0}{mL}$.
\end{solution}

\begin{solution}{exo:g10:concentration-and-dilution:5}
$C_m = c \times M = 0.10 \times 58.5 = \qty{5.85}{g/L}$, cerca de
\qty{5.9}{g/L}.
\end{solution}

\begin{solution}{exo:g10:concentration-and-dilution:6}
$V_0 = c_1 V_1 / c_0 = 0.020 \times 250.0 / 0.50 = \qty{10.0}{mL}$
(\omterm{def:g10:concentration-and-dilution:dilution}{fator de diluição} 25). Uma pipeta volumétrica de \qty{10.0}{mL} e um
balão volumétrico de \qty{250.0}{mL}.
\end{solution}

\begin{solution}{exo:g10:concentration-and-dilution:7}
O balão volumétrico tem uma única marca, feita para um volume e exata a
uma pequena fração de por cento; as graduações de um béquer são só uma
indicação grosseira. Do mesmo modo, uma pipeta volumétrica fornece um
volume com muita exatidão, enquanto uma proveta é mais larga e lida com
menos precisão.
\end{solution}

\begin{solution}{exo:g10:concentration-and-dilution:8}
Entre \qty{0.02}{mol/L} e \qty{0.04}{mol/L}: cerca de \qty{0.03}{mol/L}.
Para mais precisão, acrescentar padrões entre esses dois, por exemplo a
0.025, 0.030 e \qty{0.035}{mol/L}: uma escala mais fina onde ela é
necessária.
\end{solution}

\begin{solution}{exo:g10:concentration-and-dilution:9}
$M(\ce{C6H12O6}) = \qty{180.0}{g/mol}$; $c = 50 / 180.0 =
\qty{0.28}{mol/L}$.
\end{solution}

\begin{solution}{exo:g10:concentration-and-dilution:10}
$M = \qty{342.0}{g/mol}$, logo $n = \qty{1.00}{mol}$ e $c =
\qty{1.00}{mol/L}$. Diluída dez vezes: $c = \qty{0.100}{mol/L}$, isto é,
$C_m = 0.100 \times 342.0 = \qty{34.2}{g/L}$. Uma colherada de
\qty{0.020}{L} contém $34.2 \times 0.020 = \qty{0.68}{g}$ de sacarose.
\end{solution}

\begin{solution}{exo:g10:concentration-and-dilution:11}
(a) A mesma massa de \omterm{def:g6:solutions-and-solubility:solute}{soluto} no dobro do volume: \qty{6.0}{g/L}.
(b) O volume fica mais ou menos pela metade, a massa de \omterm{def:g6:solutions-and-solubility:solute}{soluto} não
muda: cerca de \qty{24}{g/L}.
\end{solution}

\begin{solution}{exo:g10:concentration-and-dilution:12}
\begin{enumerate}
  \item Três \omterm{def:g10:concentration-and-dilution:dilution}{diluições} de 10 vezes fazem uma \omterm{def:g10:concentration-and-dilution:dilution}{diluição} de $10^3$:
        $\qty{0.10}{mol/L} / 1000 = \qty{1.0e-4}{mol/L}$.
  \item Uma única \omterm{def:g10:concentration-and-dilution:dilution}{diluição} de 1000 vezes exigiria ou um volume minúsculo
        de \omterm{def:g10:concentration-and-dilution:dilution}{solução-estoque} (\qty{0.1}{mL} em \qty{100}{mL}), impossível
        de medir com exatidão com uma pipeta, ou um balão muito grande
        (\qty{1}{mL} em \qty{1}{L}). Três \omterm{def:g10:concentration-and-dilution:dilution}{diluições} de dez vezes usam
        pipetas e balões comuns.
\end{enumerate}
\end{solution}

\begin{solution}{exo:g10:concentration-and-dilution:13}
\begin{enumerate}
  \item $M(\ce{CaCl2}) = 40.1 + 2 \times 35.5 = \qty{111.1}{g/mol}$;
        $n = 11.1 / 111.1 = \qty{0.100}{mol}$; $c = 0.100 / 0.5000 =
        \qty{0.200}{mol/L}$.
  \item Cada \ce{CaCl2} dá um \ce{Ca^2+} e dois \ce{Cl-}:
        $[\ce{Ca^2+}] = \qty{0.200}{mol/L}$ e $[\ce{Cl-}] =
        \qty{0.400}{mol/L}$.
\end{enumerate}
\end{solution}

\begin{solution}{exo:g10:concentration-and-dilution:14}
\begin{enumerate}
  \item $r = \qty{0.60}{cm}$, $h = \qty{0.20}{cm}$:
        $V = \pi \times 0.60^2 \times 0.20 = \qty{0.23}{cm^3}$, isto é,
        \qty{0.23}{mL} de água a mais.
  \item O mesmo \omterm{def:g6:solutions-and-solubility:solute}{soluto} está em \qty{100.23}{mL} em vez de
        \qty{100.00}{mL}: a concentração fica baixa demais em
        $0.23 / 100.23 \approx
        \qty{0.23}{\percent}$.
\end{enumerate}
\end{solution}

\begin{solution}{exo:g10:concentration-and-dilution:15}
$n = 0.20 \times 0.100 + 0.10 \times 0.300 = 0.020 + 0.030 =
\qty{0.050}{mol}$ em \qty{0.400}{L}: $c = \qty{0.125}{mol/L}$. Ela não é
a média simples, \qty{0.15}{mol/L}: há três vezes mais da solução
diluída, por isso a \omterm{def:g6:pure-substances-and-mixtures:mixture}{mistura} fica mais perto de \qty{0.10}{mol/L}. É a
média ponderada pelos volumes.
\end{solution}

\begin{solution}{pb:g10:concentration-and-dilution:1}
\textbf{1.} \qty{0.9}{g} por \qty{0.100}{L}: $C_m = \qty{9.0}{g/L}$.

\textbf{2.} \qty{9.0}{g} por \qty{1000}{mL}, isto é, \qty{9.0}{mg} por
mililitro.

\textbf{3.} $m = 9.0 \times 0.500 = \qty{4.5}{g}$.

\textbf{4.} Sim: cerca de \qty{9}{g} por litro, contra uns \qty{360}{g}
por quilograma de água na saturação, umas quarenta vezes menos.

\textbf{5.} $9.0 \times 1.0 = \qty{9.0}{g}$ de sal.

\textbf{6.} $M(\ce{NaCl}) = 23.0 + 35.5 = \qty{58.5}{g/mol}$.

\textbf{7.} $c = 9.0 / 58.5 = \qty{0.154}{mol/L}$.

\textbf{8.} $n = 0.154 \times 0.500 = \qty{0.0769}{mol}$ (ou
$4.5 / 58.5$).

\textbf{9.} Cada \ce{NaCl} dá um \ce{Na+} e um \ce{Cl-}: os dois a
\qty{0.154}{mol/L}.

\textbf{10.} $0.0769 \times \num{6.02e23} = \num{4.63e22}$ \omterm{def:g9:ions:ion}{íons} sódio.

\textbf{11.} $F = 200 / 9.0 = 22.2$.

\textbf{12.} Os \qty{4.5}{g} da bolsa vêm todos da solução concentrada:
$V_0 = \qty{4.5}{g} / \qty{200}{g/L} = \qty{0.0225}{L} =
\qty{22.5}{mL}$.

\textbf{13.} $V_0 = V_1 / F = 500 / 22.2 = \qty{22.5}{mL}$.

\textbf{14.} $c_0 = 200 / 58.5 = \qty{3.42}{mol/L}$;
$c_0 V_0 = 3.42 \times 0.0225 = \qty{0.0769}{mol}$ e
$c_1 V_1 = 0.154 \times 0.500 = \qty{0.0769}{mol}$: iguais.

\textbf{15.} Cerca de $500 - 22.5 = \qty{478}{mL}$ de água estéril.

\textbf{16.} Uma pipeta graduada de \qty{25}{mL} (ou uma bureta), lida
a \qty{0.1}{mL}.

\textbf{17.} Num balão volumétrico de \qty{500.0}{mL}, completado até a
marca.

\textbf{18.} $C_m = 200 \times 23.0 / 500 = \qty{9.20}{g/L}$,
concentrada demais em $0.20 / 9.0 \approx \qty{2.2}{\percent}$.

\textbf{19.} \qty{22.5}{mL} da solução concentrada (``20 \%'') para uma
bolsa de \qty{500}{mL}.
\end{solution}
